\subsection{Выбор аппаратной схемы георадиолокатора}

Классическая радарная система является импульсной. Она излучает короткий импульс, а затем ожидает его возвращения с информацией об окружающей среде. Импульсный режим позволяет использовать одну антенну как для передачи, так и для приема: в момент передачи приемник временно блокируется, что снижает уровень самозаимствования сигнала. К преимуществам импульсных радаров также относятся простота схемотехнической реализации и отсутствие неоднозначности при интерпретации эффектов Доплера, обусловленных дальностью и скоростью \cite{noon1996stepped}.

С развитием радарных технологий для создания георадаров стали применяться более сложные аппаратные схемы, включающие радары с непрерывной частотно-модулированной волной (FMCW) и радары непрерывного излучения со ступенчатой перестройкой частоты (SFCW). На рисунке~\ref{fig:radar-waveforms-placeholder} схематически показаны временные формы сигналов, характерные для импульсных, FMCW- и SFCW-радаров.

\begin{figure}[H]
    \centering
    \fbox{\parbox[c][0.22\textheight][c]{0.82\textwidth}{\centering
    Место для схемы временных форм сигналов: FMCW, SFCW и импульсного радара.}}
    \caption{Схематический график временных форм сигналов: FMCW, SFCW и импульсного радара}
    \label{fig:radar-waveforms-placeholder}
\end{figure}

FMCW-радар использует непрерывное излучение сигнала с линейной частотной модуляцией. Передатчик и приемник разнесены, но размещаются на одной платформе. Передача и прием осуществляются одновременно, а смеситель формирует сигнал, отражающий фазовую разницу между принимаемым и опорным сигналом, пропорциональную расстоянию до цели. Таким образом, данные о расстоянии представляются в частотной области, и быстрое преобразование Фурье (FFT) позволяет получить разрешение по дальности. Преимуществами FMCW-радара являются высокая излучаемая мощность и точное разрешение по дальности. Однако метод имеет и недостатки: возможна утечка сигнала передачи в приемный тракт, значительное влияние близко расположенных объектов и ухудшение характеристик на больших дальностях.

SFCW-радар, аналогично FMCW, осуществляет передачу и прием одновременно, но в каждый момент времени излучается только одна фиксированная частота. Смеситель формирует фазовую разницу, представляющую собой комплексную постоянную, которая изменяется при переходе к следующей частоте. Таким образом, расстояние до поверхности отражения восстанавливается путем анализа совокупности фазовых и амплитудных откликов. Обладая всеми преимуществами FMCW-радара, SFCW-радар использует длительный стабильный сигнал на каждой частоте, что позволяет повысить отношение сигнал/шум (SNR) и улучшить характеристики радара.

SFCW-радары отличаются высокой гибкостью. Так, например, возможно введение задержки в начале приема каждой новой частоты, позволяющей исключить влияние ближайших целей, маскирующих слабые дальние отражения. Демодуляция и аналого-цифровое преобразование при этом значительно упрощаются. Основным ограничением SFCW-технологии является увеличенное время измерений, однако в задачах георадиолокации это, как правило, не является критичным. В силу указанных преимуществ, согласно обзорной работе \cite{sun2019advanced}, современные GPR-системы все чаще реализуются на базе принципа SFCW \cite{eide2019advanced}, в сочетании с широким набором методов цифровой обработки, направленных на повышение качества получаемых данных в конкретных прикладных сценариях. Однако в России широкополосные SFCW-системы практически не используются: в инженерной практике по-прежнему преобладают узкополосные импульсные георадары, ограниченные в глубине зондирования и разрешении.

Таким образом, GPR-системы на основе SFCW-архитектуры считаются одними из наиболее перспективных решений для задач, рассматриваемых в данной работе. По этой причине именно аппаратная реализация SFCW-георадиолокатора была выбрана для дальнейших исследований в рамках данной работы. Отсутствие серийных SFCW-георадаров отечественного производства подчеркивает актуальность исследований и разработок данной аппаратной схемы.
